% Generated draft: requires mathematical and source review.
\documentclass[11pt]{article}
\usepackage[margin=25mm]{geometry}
\usepackage{fontspec}
\setmainfont{Latin Modern Roman}
\usepackage{amsmath,amssymb,mathtools,unicode-math}
\setmathfont{Latin Modern Math}
\usepackage{enumitem,needspace,xurl,hyperref}
\setlength{\parindent}{0pt}
\setlength{\parskip}{5pt}
\setlength{\emergencystretch}{3em}
\newcommand{\N}{\mathbb{N}}
\newcommand{\Z}{\mathbb{Z}}
\newcommand{\Q}{\mathbb{Q}}
\newcommand{\R}{\mathbb{R}}
\newcommand{\C}{\mathbb{C}}
\newcommand{\F}{\mathbb{F}}
\newcommand{\A}{\mathbb{A}}
\newcommand{\PP}{\mathbb{P}}
\newcommand{\E}{\mathbb{E}}
\newcommand{\eps}{\varepsilon}
\newcommand{\dd}{\,\mathrm{d}}
\begin{document}
\section*{UnsolvedMath ID: 800421}
\textbf{Source problem code:} catalogue entry
\subsection{Definitions and mathematical statement}
Let \(q=p^r\), where \(p\) is prime and \(r\geq 1\), and let \(\F_q\) be the field with \(q\) elements. Let \(N\) denote the dimension of the projective space, and let
\[
f:\PP^N\dashrightarrow\PP^N
\]
be a dominant rational map defined over \(\F_q\). Here dominance means that the image of the domain of definition of \(f\) is Zariski dense in \(\PP^N\). Write \(\overline{\F}_q\) for an algebraic closure of \(\F_q\), and regard each \(\F_{q^k}\) as a subfield of it.

The supplied record does not define the subscript \(f\) on sets of points. For this draft, interpret it as requiring every successive forward application of \(f\) to be defined. More precisely, let \(I_f\subseteq\PP^N_{\overline{\F}_q}\) be the indeterminacy locus of \(f\). For an extension field \(\F_q\subseteq K\subseteq\overline{\F}_q\), put
\[
\PP^N(K)_f=
\left\{P\in\PP^N(K):
\begin{array}{l}
\text{there exists a sequence }(P_n)_{n\geq0}\text{ in }\PP^N(K)\\
\text{with }P_0=P,\ P_n\notin I_f,\text{ and }P_{n+1}=f(P_n)\\
\text{for every integer }n\geq0
\end{array}
\right\}.
\]
Thus these are the points whose entire forward orbit avoids indeterminacy. This interpretation of the source notation requires confirmation against the full source context.

With this interpretation, the supplied Conjecture 44 asserts that every such map satisfies both of the following statements:
\begin{enumerate}
\item[(a)] The points with well-defined forward orbit are Zariski dense:
\[
\overline{\PP^N(\overline{\F}_q)_f}^{\,\mathrm{Zar}}
=\PP^N_{\overline{\F}_q}.
\]
\item[(b)] Their proportion among points over the degree-\(k\) extensions tends to one:
\[
\lim_{\substack{k\to\infty\\k\in\Z_{\geq1}}}
\frac{\#\PP^N(\F_{q^k})_f}
{\#\PP^N(\F_{q^k})}=1.
\]
Here \(\#\) denotes cardinality; the map \(f\), the field size \(q\), and the dimension \(N\) are fixed while \(k\) tends to infinity.
\end{enumerate}
\subsection{Short English statement}
For a dominant rational self-map of projective space over a finite field, the conjecture predicts that points whose entire forward orbit is defined are Zariski dense, and that their proportion over \(\F_{q^k}\) tends to one as \(k\to\infty\). The meaning of the source's subscript \(f\) has been interpreted as avoidance of indeterminacy and requires confirmation. The registry labels this problem open; that status is supplied source data, not an independently verified assessment.
\subsection{Sources}
\begin{itemize}
\item Local UnsolvedMath registry, supplied complete source record: permanent ID \(800421\), exact source problem code \texttt{catalogue entry}, title \emph{Open Problems in Arithmetic Dynamics --- Conjecture 44}, in the AMR Open Problem Lists collection. The record attributes the problem to Benedetto, DeMarco, Ingram, Jones et al., with proposed year 2018, and labels it \emph{open}. Its supplied update date is 2026-10-16; these are registry assertions, not independently verified facts about attribution or current status.
\item Supplied scholarly reference: Benedetto, DeMarco, Ingram, Jones et al., \emph{Current Trends and Open Problems in Arithmetic Dynamics} (2018), Conjecture 44, \url{https://arxiv.org/abs/1806.04980}. The background records access on 2026-10-16 and extraction method \texttt{arxiv\_latex}. This access date and extraction information are supplied provenance; no independent source visit or verification is claimed here.
\end{itemize}
\end{document}
